\section{Faites Vos Jeux}

Zum Thema Datenschutz betrachten wir nun das Fallbeispiel \enquote{Faites Vos Jeux}\cite{kees_faites_2017} aus dem Informatik Spektrum 2017.

\subsection{AC-Games in Not}

\begin{wrapfigure}{r}{0.5\textwidth}
    
  \begin{tikzpicture}[text node/.style={font=\sffamily\small}]

    \node (walter) {\includegraphics[width=0.15\textwidth]{images/walter.png}};
    \node[text node, below=0mm of walter] (walter-name) {Walter};
    
    \node[below=2cm of walter] (acgames) {\includegraphics[width=0.15\textwidth]{images/ac_games.png}};
    \node[text node, below=0cm of acgames] (acgames-name) {AC-Games};
    
    \node[left=2.1cm of acgames] (nutzer) {\includegraphics[width=0.15\textwidth]{images/game_friends.png}};
    \node[text node, below=0cm of nutzer] (nutzer-name) {Nutzer:innen};
    
    
    \draw[-stealth] (walter-name) -- node[text node, left] (ceo) {CEO}  (acgames);
    
    % Finanzkrise 2
    
    \draw[-stealth, red] (nutzer) -- node[text node, below] (geld) {Geld} (acgames);
    
    \node[above=1cm of geld] (finanzkrise) {\includegraphics[width=0.10\textwidth]{images/money_down.png}};
    \node[text node, below=0cm of finanzkrise] (finanzkrise-name) {Finanzkrise};
    \draw[dotted, red] (finanzkrise-name) -- (geld);
    
    % App Stores machen obsolet 3
    
    \node[below=0.3cm of geld] (apps) {\includegraphics[width=0.08\textwidth]{images/smartphone.png}};
    \node[text node, below=0cm of apps] (apps-name) {App-Stores};
    \draw[dotted, red] (apps) -- (geld);
    
  \end{tikzpicture} 
  \caption{Die Situation von AC-Games.}
  \label{fig:03:acgamessituation}
\end{wrapfigure}

Walter ist seit nun fast 2 Jahrzehnten Geschäftsführer einer kleinen Spielfirma namens AC-Games, 
die ein kostenloses Spieleportal für Online-Spiele betreibt. (siehe \cref{fig:03:acgamessituation})
Lange Zeit lief das Geschäft gut, die Firma hatte einen gesunden Stamm von Nutzern und Nutzerinnen.
Die Einnahmen setzten sich aus Käufen virtueller Gegenstände, sowie Werbeeinnahmen zusammen.
Allein durch die 100 meist-zahlenden Nutzern und Nutzerinnen konnte AC-Games seine Kosten decken.

Doch dann begann vor etwa 10 Jahren die Wirtschaftskrise. 
Seitdem ging es mit AC-Games langsam bergab.
Dies war allerdings nicht allein der Wirtschaftskrise geschuldet.
Weitere Entwicklungen über die letzten 10 Jahre wirkten sich ebenfalls negativ auf AC-Games aus.

Beispielsweise wurden durch die Ausbreitung von Smartphones auch die damit einhergehenden App-Stores zum Problem für AC-Games.
Vor allem für Gelegenheitsspieler eignen sich Spiele aus solchen App-Stores 
und löst damit einen wichtigen Teil von Spieleportalen ab.
Darüber hinaus führte die Popularität von App-Stores dazu,
dass Spielentwickler weniger dazu geneigt waren, Spiele für Spieleportale zu entwickeln, im Vergleich zu App-Stores.
Darüber hinaus sanken auch sowohl die Werbeeinnahmen, sowie die Zahl der Käufe virtueller Gegenstände
und seit etwa 2-3 Jahren hat AC-Games Schwierigkeiten die eigenen Mitarbeiter und Mitarbeiterinnen zu bezahlen.

\subsection{Das Angebot}

Zu diesem Zeitpunkt bekommt Walter ein Angebot:
Ein Datenhändler namens \emph{Data Broker GmbH} möchte die Firma übernehmen und bietet eine sehr große Geldsumme.
Die hohe Geldsumme verwunderte die Mitglieder von AC-Games,
doch nach einer Diskussion entschließt sich Walter, die Firma an Data Broker zu übergeben.
Der Hauptgrund für diese Entscheidung, ist Walters Ansicht, dass AC-Games nur für Transaktionen nötige Daten speichert.
Es sei unproblematisch, diese Daten an Data Broker zu übergeben.

Doch tatsächlich speicherte AC-Games mehr Daten als von Walter angenommen.
Hierzu müssen wir an den Anfang von AC-Games zurückgehen.
Zu dieser Zeit war es Walter und seinen Mitgründern nicht sehr wichtig, auf Datenschutz zu achten.

In einem alten Blogpost stellt Walter seine Gedanken über eine sogenannte \emph{ultimative Slot-Machine} dar.
Diese soll ein Glücksspiel sein, das sich den Spielern und Spielerinnen anpasst.
Dazu müssten natürlich Verhaltensdaten gesammelt werden.
Darüber hinaus wurden weiter Überlegungen angestellt, wozu solche Daten noch dienen könnten.
Dies wurde aber ausschließlich intern besprochen.
Beispielsweise könne man Spielern künstlichem Stress aussetzen, um ihr Entscheidungsverhalten zu ihren Ungunsten zu beeinflussen.
Diese Ideen wurden allerdings wieder verworfen, inklusive der ultimativen Slot-Machine.

Allerdings wurden die angesprochenen Verhaltensdaten tatsächlich von AC-Games gesammelt und auch bis heute gespeichert.
Eine Mitarbeiterin namens Kathleen war die einzige, die einen Überblick darüber hatte, welche Daten nun tatsächlich gesammelt und gespeichert wurden.
Doch konnte sie mangels Ressourcen die Systeme nicht umprogrammieren, um die Datensammlung zu beenden.

Es wurde auch in einem Datenschutzaudit angemerkt, dass die Menge der gesammelten Daten eigenartig sei.
Allerdings wurde diesbezüglich nichts unternommen, da dies rechtlich legitim schien
und die Nutzer und Nutzerinnen in AGB eingewilligt haben, die diese Datensammlung erlauben.

Tatsächlich hat dieser Datenschutzaudit Data Broker erst auf AC-Games aufmerksam gemacht, da dieser öffentlich zugänglich war.
Im Idealfall für Data Broker müssten auch nicht die AGB geändert werden, 
sodass die Übernahme sogar ohne das Mitwissen der Nutzer und Nutzerinnen geschehen kann. 

\subsection{Der Wert der Daten}

\begin{wraptable}{l}{0.5\textwidth}
   \begin{tabular}{c|lll} 
        \textbf{Typ} & \textbf{Nutzer} \\\hline
        Stress & Bob L. & Lina M. & \dots \\\hline
        Neugier & Katie A. & Chris O. & \dots \\ \vdots & \dots & & 
    \end{tabular} 
    \caption{Hanks Kategorisierung der Nutzer und Nutzerinnen.}
    \label{tab:03:categories}
\end{wraptable}

Hank ist ein neuer Mitarbeiter bei AC-Games und der Datenbankadmin.
Auch er weiß von der Übernahme durch Data Broker und ist über die hohe Geldsumme verwundert.
Er schließt dadurch, dass die Daten von AC-Games nicht wertlos sind.
Da er von Walters alten Blog—Posts weiß und er nicht die Absicht hat, lange bei AC-Games zu bleiben,
ist er bereit etwas zu experimentieren.

Er stellt Berechnungen auf den Daten an und kategorisiert so die Nutzer und Nutzerinnen in verschiedene Gruppen. (siehe \cref{tab:03:categories})
Mithilfe dieser Daten fügt er dem In-Game Shop ein neues Element hinzu.
Er fügt einen Zähler hinzu, der die verbleibende Anzahl jedes Produkts darstellen soll.
Die angezeigte Zahl ist aber unecht.

Er passt die genaue Darstellung dieses Zählers auf die Nutzer und Nutzerinnen der Kategorien an,
die er mithilfe der Berechnungen erstellt hat.
So wird für \enquote{Stress-Typen} der Zähler zu Anfang auf eine zweistellige Zahl gesetzt und diese Zahl langsam dekrementiert.
Bei \enquote{Neugier-Typen} hingegen wird stattdessen ein \enquote{Vorrat prüfen}-Knopf angezeigt.
Wird dieser geklickt, wird ein Fenster mit dem Text \enquote{Nur noch ein Exemplar verfügbar!} angezeigt.

Am nächsten Tag prüft Hank die Verkaufszahlen und bemerkt, dass sich diese signifikant erhöht haben.
Er fragt sich, was noch mit diesen Daten angestellt werden könnte und überlegt, was er in dieser Situation tun sollte.
Er hat zwar Anweisungen bekommen, nicht tief in die Datenbank einzugreifen,
doch spielt er mit dem Gedanken, entgegen dieser Anweisungen zu handeln.
Auch entscheidet er sich dagegen, Walter von den Daten zu erzählen.

Schließlich entscheidet er sich, entgegen seinen Anweisungen, die Datenbank zu modifizieren.
Er pseudonymisiert die Nutzer-IDs und speichert eine Zuordnungstabelle auf seinem USB-Stick. 